% \documentclass[a4paper,draft,twoside,leqno]{article}
% \documentclass[a4paper,twoside,leqno]{article}

\usepackage{a4wide}

\usepackage[p,osf]{scholax}
% T1 and textcomp are loaded by package. Change that here, if you want
% load sans and typewriter packages here, if needed
\usepackage{amsmath,amsthm}% must be loaded before newtxmath
% amssymb should not be loaded
\usepackage[scaled=1.075,ncf,vvarbb]{newtxmath}% need to scale up math package
% vvarbb selects the STIX version of blackboard bold.


% The usual ams packages
% \usepackage{amssymb}
% \usepackage{amsmath}
% \usepackage{mathtools}          % for multilined environment
% \usepackage{amsthm}

% % for incrementing two counters at the same time
% \usepackage{xassoccnt}

% for the \vvvert command
% \usepackage{mathabx}

% microtype for better typesetting
\usepackage{microtype}

% for \llbracket and \rrbracket
\usepackage[only,llbracket,rrbracket]{stmaryrd}

% To produce nice calligraphic letters for Lebesgue and Hausdorff measure with
% '\mathscr{L}'
% \usepackage{mathrsfs} %% !!!!

% To insert figures.
\usepackage{graphicx}

% for mean value integral sign
% \usepackage{esint}                % load after the AMS packages

% for different enumeration styles
\usepackage{enumitem}
% \setlist[enumerate,1]{label=(\alph*),resume}
\setlist[enumerate,1]{label=(\alph*)}
\setlist[enumerate,2]{label=(\roman*)}

\usepackage[T1]{fontenc}
\usepackage[utf8x]{inputenc}
\usepackage[usenames,dvipsnames]{xcolor}

\usepackage{url}
\usepackage[hidelinks,colorlinks=true,linkcolor=blue,citecolor=darkgray,linktocpage=false]{hyperref}

% To show only referenced labels
\usepackage{mathtools}
\mathtoolsset{showonlyrefs}

% For the symbol \mathds{1} used for the characteristic function
\usepackage{dsfont}

% \usepackage{phonetic}
% \usepackage{extpfeil}
% \usepackage[all]{xy}

% to ease referencing
% \usepackage[notref,notcite]{showkeys}

% No hyphenation is better than an odd looking one
\hyphenation{Lip-schitz-ian Besicovitch}

% "empty" theorem enviroment for use in the introduction
\newtheoremstyle{citing}% name
  {3pt}%      Space above, empty = `usual value'
  {3pt}%      Space below
  {\itshape}% Body font
  {}%         Indent amount (empty = no indent, \parindent = para indent)
  {\bfseries}% Thm head font
  {.}%        Punctuation after thm head
  {.5em}%     Space after thm head: " " = normal interword space;
        %       \newline = linebreak
  {\thmnote{#3}}% Thm head spec
\theoremstyle{citing}
\newtheorem*{citing}{}

% theorem, lemma, corollary, remark, example and definition as usual
% also, intro_problem to describe the problem dealt with in the introduction
% and miniremark as environment without name to gather assumptions and the
% like
\theoremstyle{definition}
\newtheorem*{intro_problem}{Problem}

\newcounter{introcnt}
\newcounter{claimcount}
\theoremstyle{remark}
\newtheorem{claim}[claimcount]{Claim}
\newtheorem{subclaim}{Claim}[claimcount]

\swapnumbers
\theoremstyle{plain}

\newtheorem{introthm}[introcnt]{Theorem}
\newtheorem{introcol}[introcnt]{Corollary}

\newtheorem{theorem}{Theorem}[section]
\newtheorem{conjecture}[theorem]{Conjecture}
\newtheorem{lemma}[theorem]{Lemma}
\newtheorem{proposition}[theorem]{Proposition}
\newtheorem{corollary}[theorem]{Corollary}

\theoremstyle{remark}
\newtheorem{remark}[theorem]{Remark}
\newtheorem*{remark*}{Remark}
\newtheorem{example}[theorem]{Example}

\theoremstyle{definition}
\newtheorem{definition}[theorem]{Definition}
\newtheorem*{definition*}{Definition}
\newtheorem{miniremark}[theorem]{}
\newtheorem{question}[theorem]{Question}
\newtheorem*{question*}{Question}

% If equation numbers are used at all they are Roman as the enumeration is
% arabic
% \renewcommand{\theequation}{\Roman{equation}}
%\numberwithin{equation}{theorem}

% some counters
\newcounter{counter1}
\newcounter{counter2}

% Commands to conveniently subdevide proofs into several steps
\theoremstyle{definition}
\swapnumbers
\newtheorem{step}{Step} % the environment for steps
\newenvironment{proofinsteps}{\begin{proof}\setcounter{step}{0}}{\end{proof}}

% to reset equation counter each time the theorem counter increases
\usepackage{chngcntr}

% % for \lefthalfcup and \righthalfcup
% \usepackage{MnSymbol}

% \usepackage{lmodern}

% To have the same numbering for equations and enumerates
%
% Warning: putting a numbered equation inside myenumerate may spoil the
% numbering!
%
\newenvironment{myenumerate}{
  \begin{enumerate}
      \setcounter{enumi}{\value{equation}}
    }{
      \setcounter{equation}{\value{enumi}}
  \end{enumerate}
}

% Equation numbers should reset on each new item (theorem, remark, lemma, etc.).
\counterwithin*{equation}{theorem}
\counterwithin*{enumi}{theorem}


% Local Variables:
% coding: utf-8
% eval: (custom-set-variables  '(LaTeX-indent-level 4) '(LaTeX-item-indent -4) '(indent-tabs-mode nil))
% eval: (ispell-change-dictionary "british")
% eval: (flyspell-mode)
% End:
